$$s_s(x) = \begin{cases} 1 & \text{if } x \in S; \\ 0 & \text{if } x \in U + S. \end{cases}$$

The difference between B1 and B4 is subtle but important. In B4, the symbols 0 and 1 denote elements of $\mathbb{N}$ rather than $\mathbb{K}$. Of course, $c_s$ and $s_s$ are closely related, but since 1 + 1 gives a different "answer" in $\mathbb{N}$ than in $\mathbb{K}$, the characteristic function and characteristic selection are not the same thing. In particular, the correspondence $S \rightarrow s_s$ gives a natural injection of $\mathcal{P}(U)$ into $\mathbb{S}(U)$ under which $S +


There is a natural multiplication on $\mathbb{P}(U)$. Let $\mathcal{Q}, \mathcal{R} \in P(U)$ and let $\mathcal{QR}$ be the collection of nonempty subsets of the form $Q \cap R$ where $Q \in \mathcal{Q}$ and $R \in \mathcal{R}$.
